% Cote 119: offprints, reports, courses, CV and letters (typed and printed part): archive cote 119, site page 239 (PDF page 240).
% typescript. Draft compared: none.
% AI-transcribed by Claude Opus 5.5 (claude-opus-5-5, adaptive thinking, effort high), Message Batches API, batch r119_full (msgbatch_01DwZNm5Y7YsKBWyAN4sFs8w), 2026-11-13; not yet checked by eye.
\scanpage{119/239}
\mspage{- 27 -}

W. MESSING, The crystal associated to a Barsotti-%
Tate group, Lecture Notes in Math.
n° 264 (1971) Springer.

D. MUMFORD, Geometric Invariant Theory, Ergebnisse
der Math. Bd 34, 1965, Springer.

J.P. MURRE, On contravariant functors from the
category of preschemes over a field into
the category of abelian groups, Pub. Math.
n° 23, pp. 5-43 (1964).

D. QUILLEN, Homotopical Algebra, Lecture Notes in
Math. n° 43 (1967), Springer.

M. RAYNAUD, Spécialisation du Foncteur de Picard,
Publications Math. n° 38, pp. 27-76 (1970).

Mme M. RAYNAUD, Théorèmes de Lefschetz en cohomologie
cohérente et cohomologie étale, Thèse
Paris 1972 (paraîtra dans Lecture Notes
of Math.).

J.P. RAMIS et G. RUGET, Complexe dualisant et
théorèmes de dualité en géométrie analytique
complexe, Pub. Math. n° 38, pp. 77 à 91 (1970).

J. TATE, Rigid-analytic spaces, Inventiones Mathematicae,
Vol. 12, pp. 257-289.

G. TRAUTMANN, Abgeschlossenheit von Corandmoduln
und Fortsetzbarkeit kohärenter analytischer
Garben, Inventiones Math. vol. 5,
pp. 216-230 (1968).

J.L. VERDIER, J.P. RAMIS et G. RUGET, Dualité relative
en géométrie analytique complexe,
Inventiones Math., vol. 13, pp. 261-283 (1971).

\margin{239}
